% Here will go sections about the semantics of APPL, De brujin indices as well

\section{Instantiating Iris}
Theorem proving, can be split into two activities: writing definitions and proving theorems. The former doesn't seem difficult or time-consuming. It sounds reasonably
easy to come up with a skeleton for the codebase, based on a paper. After a while working on this project I have realised this is not true, and indeed this is well known
by both the theorem proving community, also true for professional mathematics and software engineering. Let's elaborate on the latter. Writing computer code is parallel
to writing definitions (of data and algorithms). So definitions are indeed important and time-consuming in Lean as well. There is no parallel to writing theorems though, programs just need to convince themselves of the correctness.

A milestone was reached when the definitions of the language and logic were written down, with a set of abstractions to convince myself this development was sound.
These definitions were composed of a separate syntax and semantics for both the language and the logic, using dependent de brujin indices for both. These definitions
would have sufficed to prove formally the correctness of all the proofs in the paper. However, of the project is not only to validate the soundness of the logic,
but also to create a usable tool for proving properties of actual probabilistic programs ergonomically (using tactics). The best way to do so is to instantiate the Iris
framework. And the interface provided by Iris was not compatible with current definitions. Below we explain why it wasn't compatible and how Iris' interface is motivated.

\subsection{Deep vs Shallow Embeddings}
When defining a language or logic, within a meta-language (lean in our case), we are presented with 2 choices. A deep embedding; inductively define the syntax,
and the denotational semantics becomes an inductive definition over the syntax. For clarity, here is an excerpt from a deep embedding of the Lilac logic:


\Todo{complete this (code from c3ca63f)}

\begin{figure}[htbp]
  \centering
  \begin{minipage}{0.9\linewidth}
\begin{leancode}
inductive Assertion : List Ty → List Ty → Type
  | top : Assertion ds rs -- ds: deterministic context, rs: random variable context
  | and  : Assertion ds rs  → Assertion ds rs → Assertion ds rs
  | sep  : Assertion ds rs  → Assertion ds rs → Assertion ds rs
  | forall (d : Ty) : Assertion (d :: ds) rs → Assertion ds rs
  | forall_rv (r : Ty) : Assertion ds (r :: rs) → Assertion ds rs
  | exists_rv (r : TyRand) : Assertion ds (r :: rs) → Assertion ds rs
  ...

def Assertion.denote {ds : List TyDet} {rs : List TyRand}
    (P : Assertion ds rs) (γ : HList (⟦·⟧) ds) (D : RV (List.TProd (⟦·⟧) rs)) (φ : PSp HC) : Prop :=
  match P with
  | top  => True
  | and P Q => P.denote γ D φ ∧ Q.denote γ D φ
  | sep P Q => ∃ φ₁ φ₂ : PSp HC, φ₁ ⋆ φ₂ ≤ some φ ∧ P.denote γ D φ₁ ∧ Q.denote γ D φ₂
  | forall d P => ∀ x : ⟦d⟧, P.denote (x :: γ) D φ
  | forall_rv r P => ∀ X : RV ⟦r⟧, P.denote γ (X ; D) φ
  | exists_rv r P => ∃ X : RV ⟦r⟧, P.denote γ (X ; D) φ
\end{leancode}
  \end{minipage}
  \caption{Excerpt of Deep Embedding of Lilac}
  \label{fig:deep-embedding}
\end{figure}


On the other hand a shallow embedding where we fuse inductive structure of the syntax and the semantic denotation function together to simply. The denotations of terms
in our object language get expressed directly in the meta language, and we lose the ability to distinguish entities in the object language from the meta-language. The previous
point, in concrete terms is that we loose the ability to do induction over the object language now. Here is an excerpt of Lilac as a shallow embedding:

\begin{leancode}

\end{leancode}

The shallow embedding had a much smaller surface and generally allowed the logic to piggyback on the forall and exists quantifiers of the meta language. 

\subsection{BI in Iris}
Remember that in \ref{sec:bi} we explained that BI (logic of bunched implications) is just abstraction which defining the standard connectives
shared by every separation logic. Here is the definition of BI as a typeclass in 
\href{https://github.com/leanprover-community/iris-lean/blob/bb24268a8a9b5ddb23224a8c1db66eadd4b9beb7/Iris/Iris/BI/BIBase.lean#L21-L36}{iris-lean} (a dependency of this project).

\begin{leancode}
/--
Require the definitions of the separation logic connectives and units on a carrier type `PROP`.
-/
class BIBase (PROP : Type u) where
  Entails : PROP → PROP → Prop
  emp : PROP
  pure : Prop → PROP
  and : PROP → PROP → PROP
  or : PROP → PROP → PROP
  imp : PROP → PROP → PROP          -- implies
  sForall : (PROP → Prop) → PROP
  sExists : (PROP → Prop) → PROP
  sep : PROP → PROP → PROP          -- separating conjunct
  wand : PROP → PROP → PROP         -- magic wand
  persistently : PROP → PROP
  later : PROP → PROP
\end{leancode}

Firstly \lean{Prop} is the universe of propositions, part of Lean's syntax. While \lean{PROP} denotes the proposition in the object logic (Lilac in our case). The
typeclass is called \lean{BIBase} since there is another typeclass \lean{BI} extending \lean{BIBase} which states the axioms of BI. As
explained in \ref{sec:bi}, the instantiation of \lean{Entails} can interpret a BI as a separation logic. \lean{Entails} says what it means for a \lean{PROP} to be true
by translating into the meta language's \lean{Prop}. \lean{persistently} and \lean{later} is omitted since the former
is an orthogonal concern and the latter is irrelevant. \lean{emp} can be understood as the unit of the monoid with which we will instantiate BI. There are several binary
operators. The most important connectives are  \lean{sForall} and \lean{sExists}. The original defintion from iris-rocq is more straightforward:

\begin{leancode}
  iForall : ∀ A : Type u', (A → PROP) → PROP
  iExists : ∀ A : Type u',  (A → PROP) → PROP
\end{leancode}

We quantify over the type A of variables that we want to quantify over, since iris can't know what type a variables a specific logic may want to quantify over, and indeed
as in Lilac there may be multiple types (Lilac has two of each quantifier). We have named the above quantifier \texttt{iForall}, with \texttt{i} short
for indexed, since when quantifying over types we refer to the types (like \texttt{A}) as \emph{type indices}. \lean{iForall} and \lean{iExists} is impredicative
(see \ref{sec:predicativity}), meaning we
quantify in universes larger than that of the type being constructed. \texttt{u'} is universe level, a natural number. Indeed, impredicativity is required here to provide complete flexibility in
what variables an instantiating logic can then quantify over.  

Now lets look at \lean{sForall} and \lean{sExists}, with \texttt{s} short for \emph{set}.  
Note that in theorem provers, Sets (say \lean{Set A}) are represented as \lean{A → Prop}. So the signature for the quantifiers says that given a set
of \lean{PROP} we construct a \lean{PROP}. The reason for this less intuitive version of quantification is that iForall (iExists) can be defined in terms of sForall
(sExists). And the quantified types are not restricted to a predetermined universe \texttt{u'}.


\subsection{First Attempt: Deeply embedded logic instantiation}
There is a result in Lilac \cite[see Appendix B.10]{lilac} called the substitution lemma which uses induction over the syntax of Lilac propositions. Supposing
that induction over the syntax of the logic is important/necessary, a next step is to try to instantiate \lean{BIBase} above with the deep embedding from Fig 
\ref{fig:deep-embedding}. We may try something like so:

\begin{leancode}
instance (ds rs: List Ty) : BIBase (Assertion ds rs) where
  Entails P Q    := ∀ γ D φ, P.denote γ D φ → Q.denote γ D φ
  and            := Assertion.and
  ...
\end{leancode}

Where we simply instantiate each of the constructors with syntax, and called the denotation function in \lean{Entails}. This works nicely until we get to the instantiation
of the quantifiers. We now need some syntactic version of \lean{sForall} and not just the 
\lean{forall}/\lean{forall_rv} from Fig \ref{fig:deep-embedding}. Though the latter is all we need for our logic instantiating
iris demands the flexibility to quantify over any type, not just the object language types. Let's introduce the following syntax into our
language:

\begin{leancode}
inductive Assertion' : List TyDet → List TyRand → Type
  | sForall : (Assertion' ds rs → Prop) → Assertion' ds rs
  ...
\end{leancode}

This looks like a valid lean definition, but it is not, since it fails a syntactic requirement of inductive types called \emph{strict positivity}.
The type being defined inductively may only appear in \emph{postiive} locations, which means they are not allowed to occur in the input of an arrow
type. Even if we suppose the \lean{iForall} definition was used in BI, we still end up with nuanced problems with mixing semantics into the syntax
because of the quantifiers. We spare further details, but the conclusion is that the impredicativity of quantifiers in \lean{BI}
disallows usage of a deep embedding of the logic.

\subsection{Shallow embedding of the logic}


\section{Deparametrisation of LProp}

A major concern throughout the mechanisation was the definition of \texttt{LProp}, and how amenable it might be to variable substitution. Remember that we were forced to forego
the convenience of an induction principle over the lilac connectives (or constructors in this context). This was indeed due to the generality of iris without which we wouldn't
be able to use it in the first place.
This seemingly insignificant issue tended to
have a significant effect on the decisions made. We explain below,
why the notion of substitution as defined in the paper, is incompatible with iris. % even though iris doesn't have particularly strong requirements on this matter. 

% Then present parts of the paper, that were
% not followed in this mechanisation, and the justification for the definition of LProp chosen after this deliberation.

Lilac has two contexts, the random and deterministic contexts. The latter is only required when the logic supports reasoning about conditioning. Remember that conditioning takes random variable
and fixes its values, i.e. makes it deterministic. And only for that deterministic value do we need the deterministic context and otherwise, every variable in our language corresponds to a random variable.
So we will focus on the random contexts. Take the quantifier $\forall_{rv}$:

\[
\gamma, D, \mathcal{P} \vDash \forall_{rv} X : A.\ P \qquad iff \qquad \text{for all } X : RV A (\gamma, (D, X), \mathcal{P}) \vDash P 
\]

We are quantifying over $RV A$, which means variables in Lilac are defined to have this type. Remember $RV A$ means $HC \xrightarrow{m} A$ (with additional finite footprint property not relevant here).
Substitution is the process of transforming from one LProp to another by specifying a mapping from one variable to another. However, definition of substitution given in the paper looks like this:

\[
S \in \llbracket \Delta' \rrbracket \to \llbracket \Delta\rrbracket 
\]

\[
\begin{aligned}
\top[S] &= \top \\
\bot[S] &= \bot \\
(P \wedge Q)[S] &= (P[S] \wedge Q[S]) \\
(P \vee Q)[S] &= (P[S] \vee Q[S]) \\
(P \to Q)[S] &= (P[S] \to Q[S]) \\
(P * Q)[S] &= (P[S] * Q[S]) \\
(P -* Q)[S] &= (P[S] \wand Q[S]) \\
(\forall_{\mathrm{rv}} X:A.\,P)[S] &= \forall_{\mathrm{rv}} X:A.\,P[S \times 1_A] \\
(\exists_{\mathrm{rv}} X:A.\,P)[S] &= \exists_{\mathrm{rv}} X:A.\,P[S \times 1_A] \\
...
\end{aligned}
\]

It is helpful to have the following mental picture of the meaning of an LProp as a function $HC \xrightarrow{m} \llbracket \Delta \rrbracket \xrightarrow{m} A$, pick an element $\omega$ of HC
(which is like sampling a random number). Once we do this the randomness is gone and the maps forward are deterministic. $\omega$ maps to values of variables, $\llbracket \Delta \rrbracket$,
as an intermediate step. Then these variables deterministically map to some expression of type A. With this in mind note that there's a mismatch in what we would like substitution to be, replacing
$X : HC \xrightarrow{m} A$ with $Y : HC \xrightarrow{m} A$ which is used in an LProp, with the above definition of substitution, which only starts at the intermediate $\llbracket \Delta \rrbracket$ step.

At a high level the Iris proof mode (indeed any proof mode) works by splitting an LProp into a context of premises and variables, and a goal. The variables are of type $RV A$ and if the substitutions are
not with $RV A$ as domain, the types just won't line up.

% To make things more concrete observe the substitution rule above for $\forall_{rv}$. Thi
% Let's give a more concrete reason for why this is problematic. The iris proof mode proves an LProp by splitting it into the context which contains premises and universally quantified
% variables and the goal. The user tries maintains a contact of premises and universally quantified variables When using the iris proof mode we generally attempt to bring ``intro'' (using the intro tactic) 

% entails fixing an arbitrary value for that variable and 
% from the random context, converts the random variable to a deterministic

In order to explain further we need an example of where substitution needs to be used. One canonical usage of substitution is in the proof rule \texttt{wp\_ret}:
\[
Q[[M]/X] \vdash \mathrm{wp}([[ret M]], X.Q)
\]

saying that the random variable obtained by sampling the program \( \llbracket ret M \rrbracket \) is just always the denotation of M, \( \llbracket M \rrbracket \). \( \llbracket M \rrbracket \) substitutes the
variables X which is free in the LProp $Q$, so substitution is essentially removal of free variables. Initial attempts at defining \( wp\_ret \) tried to adapt the proposed at the \( \llbracket \Delta \rrbracket \)
level, but this led to a painfully convoluted system before finally hitting a dead end with this approach. We try to outline the old design with the old definitions of LProp. At a high level, the above
notion of substitution can be made to work 

The development got quite far with the \lean{LProp} being parametrised with the type of the deterministic context. However, that started causing
serious issues when stating the proof rules, because of the book-keeping that needs to be done with always keeping track of the number of variables
in context. Not to mention this bookkeeping is quite difficult since we always need to keep proving that the the contexts are a measurableSpace.


\subsection{Substitution inconsistency}
Lilac
